\chapter{Bayesowskie testowanie hipotez, iloraz Bayesa i paradoksy}
\label{ch:bayes-model-testing}
\paragraph{Powiązania i odsyłacze.} Porównaj definicję klasycznej wartości p w rozdziale~\ref{ch:testing} i decyzje bayesowskie w rozdziale~\ref{ch:bayes-estimation-prediction}.


\section{Prawdopodobieństwo hipotezy a zgodność danych z hipotezą}
W testowaniu klasycznym z rozdziału~\ref{ch:testing}
obliczamy prawdopodobieństwo odpowiednio skrajnych danych
\emph{pod warunkiem} hipotezy zerowej. Bayesowskie podejście
porównuje prawdopodobieństwa samych modeli \emph{po danych},
co wymaga zadania rozkładów a priori również na \emph{modele}.
Są to inne wielkości. W szczególności wartość p nie jest
prawdopodobieństwem prawdziwości $H_0$, a iloraz Bayesa nie
jest bez dodatkowych założeń prawdopodobieństwem błędu.

\begin{definicja}[Porównanie dwóch modeli i iloraz Bayesa]
\label{def:bayes-factor}
Niech $\mathcal M_0$ i $\mathcal M_1$ będą dwoma modelami
obserwacji $X$ w tej samej przestrzeni mierzalnej, a
$w_0,w_1>0$, $w_0+w_1=1$, ich wcześniejszymi
prawdopodobieństwami. W modelu $j\in\{0,1\}$
niech $\Theta_j$ będzie przestrzenią parametrów,
$\pi_j$ \emph{właściwym} priorem i
$p_j(x\mid\theta_j)$ gęstością lub funkcją masy
względem tej samej, ustalonej miary $\nu$ na przestrzeni
danych. Definiujemy \emph{zintegrowaną wiarygodność}
$m_j(x)=\int_{\Theta_j}p_j(x\mid\theta_j)\,\dd\pi_j(\theta_j)$.
Jeżeli $0<m_j(x)<\infty$ dla obu $j$, to
$BF_{10}(x)=m_1(x)/m_0(x)$ nazywamy \emph{ilorazem
Bayesa} na korzyść modelu $\mathcal M_1$. Analogicznie
$BF_{01}=1/BF_{10}$.
\end{definicja}

\begin{twierdzenie}[Posterior odds = prior odds razy Bayes factor]
\label{thm:bayes-odds}
W założeniach definicji~\ref{def:bayes-factor}, dla danych
$x$ z $0<w_0m_0(x)+w_1m_1(x)<\infty$ mamy
\[
\frac{\Prob(\mathcal M_1\mid X=x)}
     {\Prob(\mathcal M_0\mid X=x)}
=BF_{10}(x)\,\frac{w_1}{w_0}.
\]
W szczególności $\Prob(\mathcal M_1\mid x)
=w_1m_1(x)/(w_0m_0(x)+w_1m_1(x))$.
\end{twierdzenie}
\begin{proof}
Zmienna indeksująca model ma rozkład dyskretny
$\Prob(J=j)=w_j$, a warunkowa gęstość danych przy $J=j$
to $m_j(x)$ po scałkowaniu parametru $\theta_j$.
Wzór Bayesa dla tej dwupunktowej przestrzeni
parametrów daje $\Prob(J=j\mid x)=w_jm_j(x)/
(w_0m_0(x)+w_1m_1(x))$. Podzielenie odpowiednich
tożsamości kończy dowód.
\end{proof}

\begin{uwaga}[Znaczenie doboru priorów]
Iloraz Bayesa zależy od \emph{rozkładów parametrów wewnątrz
modeli}, nie tylko od ich najlepszych dopasowań. Zastąpienie
$\pi_1$ rozkładem znacznie szerszym może zmienić
$BF_{10}$ nawet wtedy, gdy maksimum wiarygodności i
obserwowane dane pozostają takie same. Priory niewłaściwe
nie pozwalają na zwykłe porównanie modeli z niezależnymi
nieustalonymi stałymi normalizacji
(por. rozdział~\ref{ch:bayes-foundations}).
\end{uwaga}

\section{Dokładny test punktowej hipotezy o odsetku wad}
\begin{twierdzenie}[Iloraz Bayesa: punktowy Bernoulli kontra beta]
\label{thm:bernoulli-point-bf}
Niech $n\ge1$ i $X_1,\ldots,X_n\in\{0,1\}$. Porównujemy
$\mathcal M_0: X_i\iid\operatorname{Bernoulli}(p_0)$,
gdzie $0<p_0<1$, z
$\mathcal M_1: p\sim\operatorname{Beta}(a,b)$,
$a,b>0$, i $X_i\mid p\iid\operatorname{Bernoulli}(p)$.
Jeżeli $S=\sum_iX_i=s$, to
\[
BF_{10}(s)=\frac{B(a+s,b+n-s)}{B(a,b)\,p_0^s(1-p_0)^{n-s}}.
\]
Ten sam wzór otrzymujemy z obserwacji pełnego wektora
lub wyłącznie sumy $S$.
\end{twierdzenie}
\begin{proof}
Dla pełnej próbki $x=(x_1,\ldots,x_n)$ jej funkcja masy
w modelu zerowym wynosi
$p_0^s(1-p_0)^{n-s}$, a w alternatywnym całka
$\int_0^1p^s(1-p)^{n-s}\,\dd\operatorname{Beta}(a,b)(p)$.
Całka jest równa $B(a+s,b+n-s)/B(a,b)$.
Gdy obserwujemy tylko sumę $S$, obie funkcje masy mają
dodatkowy wspólny czynnik $\binom ns$, który skraca się
w ilorazie. W tym modelu $S$ jest statystyką dostateczną,
a zatem agregacja nie gubi informacji o $p$.
\end{proof}

\begin{przyklad}[Kontrola jakości: siedem wad na sto]
Testujemy $\mathcal M_0:p=0{,}10$ wobec alternatywy
$\mathcal M_1:p\sim\operatorname{Beta}(1,1)$, używając
$S=7$, $n=100$. Z powyższego wzoru
$BF_{10}\approx0{,}11138$, a $BF_{01}\approx8{,}9784$.
Przy równych priorowych prawdopodobieństwach modeli
$\Prob(\mathcal M_0\mid S=7)\approx0{,}8998$.
Nie wolno jednak opisać tego jako prawdopodobieństwa, że
$H_0$ jest prawdziwa \emph{bez podania priorycznego
prawdopodobieństwa samej hipotezy punktowej}.
Dla alternatywy $\operatorname{Beta}(2,18)$ otrzymujemy
$BF_{10}\approx0{,}70745$, czyli znacznie inny wynik,
choć obserwacje i $p_0$ się nie zmieniły.
Przy alternatywie $\operatorname{Beta}(20,80)$
$BF_{10}\approx0{,}03371$. Widać stąd konieczność
analizy wrażliwości na prior alternatywy.
\end{przyklad}

\begin{uwaga}[Hipoteza punktowa potrzebuje dodatniej masy modelowej]
Jeżeli jedynym priorem na ciągły parametr $p$ jest rozkład
beta, to $\Prob(p=p_0)=0$ przed obserwacją i po niej,
niezależnie od danych o dodatniej wiarygodności.
Aby mówić o dodatnim posteriorowym prawdopodobieństwie
punktowej hipotezy $p=p_0$, trzeba wprowadzić model
przypisujący jej dodatnią masę, np. prior mieszany
$w_0\delta_{p_0}+w_1\operatorname{Beta}(a,b)$.
Takie porównanie modeli nazywa się niekiedy priorem
\emph{spike-and-slab}. Nie należy mylić z nim posteriorowego
prawdopodobieństwa nierówności $p>p_0$, które możemy obliczać
bez atomu w $p_0$.
\end{uwaga}

\section{Decyzja oparta na posteriorowych prawdopodobieństwach modeli}
\begin{twierdzenie}[Próg decyzji przy nierównych kosztach błędów]
\label{thm:bayes-hypothesis-decision}
Załóżmy model z twierdzenia~\ref{thm:bayes-odds}.
Wybieramy pomiędzy działaniami $a_0$ (uznajemy model
$\mathcal M_0$) i $a_1$ (uznajemy $\mathcal M_1$).
Strata wynosi zero przy poprawnym wyborze,
$C_{10}>0$ przy wyborze $a_1$, gdy prawdziwy jest
$\mathcal M_0$, i $C_{01}>0$ przy wyborze $a_0$,
gdy prawdziwy jest $\mathcal M_1$. Wtedy $a_1$
ściśle minimalizuje stratę posteriorową wtedy i tylko wtedy, gdy
\[
 BF_{10}(x)>\frac{C_{10}}{C_{01}}\frac{w_0}{w_1}.
\]
Przy równości obie decyzje są optymalne.
\end{twierdzenie}
\begin{proof}
Posteriorowa strata $a_1$ równa się
$C_{10}\Prob(\mathcal M_0\mid x)$, a strata $a_0$ wynosi
$C_{01}\Prob(\mathcal M_1\mid x)$.
Porównanie tych liczb i podstawienie ilorazu z
 twierdzenia~\ref{thm:bayes-odds} daje tezę.
\end{proof}

\begin{przyklad}[Jak koszt błędnej decyzji zmienia próg?]
Niech $w_0=w_1=1/2$. Jeśli koszt fałszywego odrzucenia
$\mathcal M_0$ wynosi $10$ jednostek, a koszt błędnego
zaakceptowania $\mathcal M_0$ wynosi $2$, to wybieramy
$\mathcal M_1$ dopiero przy $BF_{10}>5$.
Żaden uniwersalny próg, np. $BF_{10}>3$ lub $BF_{10}>10$,
nie zastąpi rzeczywistej analizy kosztów.
\end{przyklad}

\section{Paradoks Jeffreysa--Lindleya: to samo z, różny wniosek}
\begin{twierdzenie}[Iloraz Bayesa w normalnym teście średniej]
\label{thm:jeffreys-lindley}
Niech $n\ge1$ oraz $X_1,\ldots,X_n\mid\mu$ będą
niezależne o rozkładzie $\mathcal N(\mu,\sigma^2)$
ze znanym $\sigma>0$. Porównujemy
$\mathcal M_0:\mu=0$ z
$\mathcal M_1:\mu\sim\mathcal N(0,\tau^2)$,
$\tau>0$, przy tym samym rozkładzie warunkowym
obserwacji. Niech $z=\sqrt n\,\overline x/\sigma$
i $r=n\tau^2/\sigma^2$. Wtedy
\[
 BF_{01}(x)=\sqrt{1+r}\,
 \exp\!\left(-\frac{r}{1+r}\frac{z^2}{2}\right).
\]
Dla ustalonego $z\in\R$, $\sigma>0$, $\tau>0$ i $n\to\infty$
wartość $BF_{01}(x)$ dąży do nieskończoności, chociaż
klasyczna dwustronna wartość p równa
$2(1-\Phi(|z|))$ pozostaje stała.
\end{twierdzenie}
\begin{proof}
Przy znanej wariancji $\sigma^2$ resztowa suma kwadratów
$\sum_i(x_i-\overline x)^2$ ma ten sam czynnik gęstości
w obu porównywanych modelach i skraca się w ilorazie.
Wystarczy zatem porównać gęstości $\overline X$.
W modelu zerowym $\overline X\sim\mathcal N(0,\sigma^2/n)$.
W alternatywnym $\overline X=\mu+(\overline X-\mu)$,
więc po scałkowaniu normalnego prioru mamy
$\overline X\sim\mathcal N(0,\sigma^2/n+\tau^2)$.
Iloraz gęstości dwóch rozkładów normalnych w punkcie
$\overline x$ daje podany wzór.
Gdy $n\to\infty$ przy ustalonym $z$, to $r\to\infty$,
czynnik wykładniczy zbiega do $e^{-z^2/2}>0$,
a pierwiastek $\sqrt{1+r}\to\infty$. Natomiast rozkład
statystyki $Z=\sqrt n\,\overline X/\sigma$ pod $\mathcal M_0$
jest standardowo normalny dla każdego $n$, więc
wartość p zależy wyłącznie od $|z|$.
\end{proof}

\begin{przyklad}[Istotność 5\% i coraz większa przewaga modelu zerowego]
Przyjmijmy $\sigma=\tau=1$ i zaobserwowane $z=2$.
Dwustronna klasyczna wartość p to
$2(1-\Phi(2))\approx0{,}04550$ dla każdego $n$.
Jednak przy $n=100$, $1000$ i $10000$ dostajemy
odpowiednio $BF_{01}\approx1{,}3873$, $4{,}2904$ oraz
$13{,}5369$. Przy priorowych szansach $1:1$
posteriorowe prawdopodobieństwa modelu zerowego wynoszą
odpowiednio około $0{,}581$, $0{,}811$ i $0{,}931$.
To nie jest sprzeczność rachunkowa. Test klasyczny i iloraz
Bayesa odpowiadają na różne pytania, a przy wzroście $n$
zachowujemy tę samą \emph{wartość statystyki z}, czyli
obserwowana średnia maleje jak $n^{-1/2}$.
Nie należy mylić tej konstrukcji z zachowaniem
stałej, niezerowej wartości prawdziwego parametru $\mu$.
\end{przyklad}

\begin{center}
\begin{minipage}{.82\textwidth}
\centering
\begin{tikzpicture}[>=Latex,scale=.88,every node/.style={font=\small}]
\draw[->,thick] (0,0)--(6.7,0) node[right] {$n$};
\draw[->,thick] (0,0)--(0,5.15) node[above] {$BF_{01}$};
\foreach \x/\lab in {1/100,3/1000,5/10000} {
  \draw (\x,.08)--(\x,-.08) node[below=3pt] {\lab};
}
\foreach \y/\lab in {0/0,1.5/5,3/10,4.5/15} {
  \draw (.08,\y)--(-.08,\y) node[left=3pt] {\lab};
}
\draw[very thick,blue!60!black] (1,.4162)--(3,1.2871)--(5,4.0611);
\foreach \x/\y in {1/.4162,3/1.2871,5/4.0611} {
  \fill[blue!60!black] (\x,\y) circle (2.4pt);
}
\node[anchor=west,font=\footnotesize,fill=white,inner sep=2pt] at (0.45,4.65)
 {$z=2$, $\sigma=\tau=1$};
\end{tikzpicture}

{\small\emph{Rysunek 20.1.} Przy stałej statystyce $z=2$
wartość p pozostaje równa około $0{,}0455$, podczas gdy
$BF_{01}$ rośnie. Punkty oznaczają $n=100,1000,10000$;
połączono je tylko dla czytelności.}
\end{minipage}
\end{center}

\begin{uwaga}[Zbyt szeroka alternatywa — paradoks Bartletta]
W twierdzeniu~\ref{thm:jeffreys-lindley} także przy
ustalonym $n$ i zaobserwowanym $z$ zachodzi
$BF_{01}\to\infty$ dla $\tau\to\infty$.
Zwiększenie rozproszenia prioru alternatywy obniża jej
zintegrowaną wiarygodność w konkretnym punkcie danych,
choć nie zmienia wartości maksymalnej wiarygodności.
To ważny argument przeciw bezrefleksyjnemu określaniu
bardzo szerokiego prioru jako „neutralnego” w porównywaniu
modeli. Granica $\tau\to\infty$ nie uprawnia do zastąpienia
prioru właściwego priorem niewłaściwym w ilorazie Bayesa.
\end{uwaga}

\section{Sekwencyjne testowanie: kiedy wolno patrzeć na dane wielokrotnie?}
W klasycznym testowaniu bez korekty wielokrotne sprawdzanie,
czy wartość p przekroczyła próg, może znacząco zwiększać
częstość błędu pierwszego rodzaju. W specjalnej sytuacji
\emph{prostej} hipotezy zerowej iloraz Bayesa ma przydatną
własność martyngałową, lecz należy precyzyjnie sprawdzić
warunki, zanim wykorzystamy ją do zatrzymywania eksperymentu.

\begin{twierdzenie}[Iloraz Bayesa jako nieujemny martyngał]
\label{thm:bayes-factor-martingale}
Niech $(\mathsf S,\mathcal S,\nu)$ będzie przestrzenią
z miarą $\sigma$-skończoną. Przy $\mathcal M_0$ zmienne
$X_1,X_2,\ldots$ są niezależne z tą samą dodatnią
$\nu$-prawie wszędzie gęstością $f_0$.
Przy alternatywie $\mathcal M_1$ ustalamy właściwy prior
$\pi$ na przestrzeni parametrów $\Theta_1$ i
wspólnie mierzalne gęstości $f_\theta$ względem $\nu$;
przy znanym $\theta$ obserwacje są niezależne o gęstości
$f_\theta$. Niech $\mathcal F_n=\sigma(X_1,\ldots,X_n)$,
$\mathcal F_0=\{\varnothing,\Omega\}$, oraz
\[
 B_n=\int_{\Theta_1}\prod_{i=1}^n
    \frac{f_\theta(X_i)}{f_0(X_i)}\,\dd\pi(\theta),
 \qquad B_0=1.
\]
Wtedy przy $\mathcal M_0$ proces $(B_n)_{n\ge0}$ jest
nieujemnym martyngałem, $\E_0 B_n=1$, oraz dla
każdego $\alpha\in(0,1)$ zachodzi
\[
 \Prob_0\!\left(\sup_{n\ge0}B_n\ge\frac1\alpha\right)
 \le\alpha.
\]
\end{twierdzenie}
\begin{proof}
Dla ustalonego $\theta$ wyrażenie
$L_n(\theta)=\prod_{i=1}^n f_\theta(X_i)/f_0(X_i)$
ma przy $\mathcal M_0$ wartość oczekiwaną jeden, ponieważ
$\int_{\mathsf S} f_\theta\,\dd\nu=1$.
Z niezależności nowej obserwacji wynika
$\E_0[L_{n+1}(\theta)\mid\mathcal F_n]=L_n(\theta)$.
Tonelli i warunkowa wersja tego twierdzenia pozwalają
zamienić całkę po $\pi$ z warunkową wartością oczekiwaną.
Otrzymujemy $\E_0[B_{n+1}\mid\mathcal F_n]=B_n$,
co wraz z $\E_0B_n=1$ daje pierwszą tezę.
Dla $a>0$ i $N\in\N$ zatrzymajmy proces w
$\tau=\inf\{n\ge0:B_n\ge a\}\wedge N$.
Twierdzenie o zatrzymywaniu dla ograniczonego czasu daje
$1=\E_0 B_\tau\ge a\Prob_0(\max_{0\le k\le N}B_k\ge a)$.
Po podstawieniu $a=1/\alpha$ i przejściu $N\to\infty$
przez ciągłość prawdopodobieństwa dla rosnących zdarzeń
otrzymujemy tezę.
\end{proof}

\begin{uwaga}[Dokładna granica zastosowania]
Próg $BF_{10}\ge1/\alpha$ zapewnia w powyższym
\emph{prostym} modelu zerowym kontrolę prawdopodobieństwa
fałszywego przekroczenia poziomu $\alpha$, nawet przy
kolejnym sprawdzaniu danych. Nie jest to twierdzenie
o dowolnej złożonej hipotezie zerowej, priorycznych
prawdopodobieństwach modeli ani o poziomie klasycznego
testu opartego na zwykłej wartości p. Także zmiana prioru
alternatywy w reakcji na bieżące dane narusza przyjętą
konstrukcję martyngału. Do zastosowań konieczne są
wstępnie ustalone modele i reguły raportowania.
\end{uwaga}

\begin{przyklad}[Adaptacyjny prior niszczy gwarancję sekwencyjną]
\label{ex:adapt-bf-fails}
Niech pod hipotezą zerową $X\sim\operatorname{Bernoulli}(1/2)$.
Rozpatrzmy dwie dopuszczalne alternatywy punktowe:
$P_0(X=0)=1$ oraz $P_1(X=1)=1$. Jeśli \emph{przed}
obserwacją ustalimy jedną z nich, iloraz wiarygodności
$L(X)=P_j(\{X\})/(1/2)$ ma średnią jeden pod hipotezą
zerową. Jeśli jednak \emph{po} zobaczeniu $X$ wybierzemy
alternatywę $P_X$, to otrzymamy $L(X)=2$ przy każdym wyniku,
a zatem $\E_0 L(X)=2$, nie zaś jeden.
Przy progu $1/\alpha=2$ (czyli $\alpha=1/2$) fałszywe
przekroczenie występuje z prawdopodobieństwem jeden.
Jest to konkretny kontrprzykład pokazujący, dlaczego prior
alternatywy w twierdzeniu~\ref{thm:bayes-factor-martingale}
musi być zadany przed obserwacją, a nie dobrany do danych.
\end{przyklad}

\section{Jak w praktyce wybierać między testem klasycznym a Bayes factor?}
\begin{center}\small
\begin{tabular}{>{\raggedright\arraybackslash}p{.27\textwidth}>{\raggedright\arraybackslash}p{.31\textwidth}>{\raggedright\arraybackslash}p{.27\textwidth}}
\toprule
\textbf{Cel analizy} & \textbf{Narzędzie} & \textbf{Interpretacja i ograniczenie}\\
\midrule
Kontrola błędu I rodzaju przy ustalonym $H_0$ & Test klasyczny i wartość p & Poziom obowiązuje przy warunkach testu.\\
Porównanie dwóch opisanych modeli & $BF_{10}$ i posterior modelowy & Trzeba podać oba priory i szanse modeli.\\
Ocena $\Prob(p>p_0\mid x)$ & Posterior ciągły i ogon rozkładu & Nie jest to $BF$ hipotezy punktowej.\\
Decyzja o koszcie różnych błędów & Posterior odds i funkcja straty & Próg zależy od rzeczywistych kosztów.\\
Wielokrotne sprawdzanie w czasie & Z góry ustalony proces ilorazu wiarygodności & Kontrola przekroczenia dla prostej $H_0$; nie dowolny BF.\\
Słaby prior alternatywy & Analiza wrażliwości $BF$ & Zbyt szeroki prior może sprzyjać $H_0$.\\
\bottomrule
\end{tabular}
\end{center}

\Needspace{23\baselineskip}
\begin{lstlisting}[language=Python,caption={Kompletny program: iloraz Bayesa dla wadliwych produktów i paradoks Jeffreysa--Lindleya (Python, SciPy).}]
from math import exp, log, log1p, sqrt
from scipy.special import betaln
from scipy.stats import norm

n, s, p0 = 100, 7, 0.10
for a, b in [(1, 1), (2, 18), (20, 80)]:
    log_bf10 = (betaln(a+s, b+n-s) - betaln(a, b)
                - s*log(p0) - (n-s)*log1p(-p0))
    bf10 = exp(log_bf10)
    post_prob_h0 = 1.0 / (1.0 + bf10)  # equal prior model odds
    print("Alternative Beta", (a, b),
          "BF10=", round(bf10, 5),
          "P(H0|data)=", round(post_prob_h0, 5))

sigma, tau, z = 1.0, 1.0, 2.0
print("Two-sided classical p-value:", 2*norm.sf(abs(z)))
for size in [100, 1000, 10000]:
    r = size * tau**2 / sigma**2
    bf01 = sqrt(1+r) * exp(-0.5*z*z*r/(1+r))
    print("n=", size, "BF01=", round(bf01, 5),
          "P(H0|data)=", round(bf01/(1+bf01), 5))
\end{lstlisting}

\begin{zadanie}[Kiedy alarm jest optymalny?]
Dwa modele mają priorowe prawdopodobieństwa $w_0=0{,}8$
i $w_1=0{,}2$. Otrzymano $BF_{10}=4$.
Fałszywe odrzucenie modelu $\mathcal M_0$ kosztuje $6$,
a przeoczenie modelu $\mathcal M_1$ kosztuje $10$.
Które działanie minimalizuje posteriorową stratę?
\par\emph{Rozwiązanie.} Posteriorowe szanse modelu
$\mathcal M_1$ wynoszą $4\cdot(0{,}2/0{,}8)=1$,
zatem oba modele mają posteriorowe prawdopodobieństwa po $1/2$.
Ryzyko działania $a_1$ wynosi $6\cdot1/2=3$,
a działania $a_0$ wynosi $10\cdot1/2=5$.
Optymalny jest wybór $a_1$. Równoważnie sprawdzamy
$BF_{10}>(6/10)(0{,}8/0{,}2)=2{,}4$.
\end{zadanie}

% ================================================================
% ROZDZIAL 21 -- MODELE HIERARCHICZNE
% ================================================================
